\documentclass[10pt,a4paper]{amsart}
\usepackage[margin=19mm]{geometry}
\usepackage{amsmath,amssymb,mathtools}
\usepackage[scr=rsfs,cal=euler]{mathalfa}
\usepackage{xfrac}
\usepackage[unicode,hidelinks,bookmarks=false]{hyperref}
\newcommand{\ie}{i.e.\ }
\newcommand{\cf}{cf.\ }
\newcommand{\resp}{respectively\ }
\newcommand{\Subsection}{\subsection}
\providecommand{\nobreakdash}{\nobreak\hbox{-}}
\setlength{\emergencystretch}{2em}
\multlinegap0pt
\usepackage{bm}
\usepackage{bbm}
\usepackage{dsfont}
\usepackage{xspace}
\usepackage{cancel}
\usepackage{mathtools}
\usepackage{amsthm}
\makeatletter
\@ifundefined{theorem}{\newtheorem{theorem}{Theorem}}{}
\@ifundefined{lemma}{\newtheorem{lemma}[theorem]{Lemma}}{}
\makeatother
\makeatletter
\DeclareMathOperator{\Ext}{Ext}
\DeclareMathOperator{\Gdim}{G-dim}
\DeclareMathOperator{\Hom}{Hom}
\DeclareMathOperator{\RHom}{\mathsf{RHom}}
\expandafter\ifx\csname bfchunk\endcsname\relax
\newenvironment{bfchunk}{\begin{chunk}\textbf}{\end{chunk}}
\fi
\makeatother
\title[Ascent and descent of Gorenstein homological properties]{Ascent and descent of Gorenstein homological properties}
\author{Jian Liu, Wei Ren}
\date{Source: arXiv:2309.02372v4 (CC BY 4.0)}
\begin{document}
\maketitle

\noindent\emph{Source and licence.} Ascent and descent of Gorenstein homological properties. Version arXiv:2309.02372v4 (CC BY 4.0), distributed under the Creative Commons Attribution 4.0 International licence, as recorded in the arXiv licence metadata for this exact version and shown on the abstract page https://arxiv.org/abs/2309.02372v4. Full rights notice: licenses/arxiv-2309.02372-CC-BY-4.0.txt.\par\smallskip

\noindent\textit{Editorial scope.} Counted unit: the Theorem whose source label is \texttt{ADTF} in the article (source0.tex lines 718--727, source label \texttt{ADTF}), delivered together with the article's own complete original proof of it (source0.tex lines 729--756, one \texttt{proof} environment of 28 lines). No mathematics is written, repaired or re-derived: statement and proof are the article's own text line for line. Required context retained uncounted: observation (lines 596--604); TFtest (lines 707--709). Every definition, display and label of those lines is retained unchanged. Omitted from this package and not redistributed: 1--595, 605--706, 710--717, 728--728, 757--996. Nothing inside a retained excerpt is omitted or altered: every retained body is delivered exactly as the article's lines are written, so no omission receipt of any kind accompanies this package. Citations: the retained excerpts carry no citation command at all, so no bibliography entry of the article is transcribed and no reference is redistributed. Reference adaptation: none was needed and none was made; every reference inside the delivered unit resolves inside it, so no unresolved reference or citation is produced. Printed slips kept literal: no printed word, symbol, hypothesis or number of the counted statement or of its proof was altered. Layout and class: the article's own class is \texttt{amsart} with packages including amsfonts, amssymb, amsmath, enumerate, verbatim, mathtools, tikz, bm, mathrsfs, tikz-cd, hyperref, comment, stmaryrd, amsthm; the wrapper is the lane's pinned \texttt{amsart} editorial wrapper at 10pt on A4, which restates the theorem-like environments the retained text needs and adds the article's own macro definitions that the retained text needs (\texttt{\textbackslash Ext}, \texttt{\textbackslash Gdim}, \texttt{\textbackslash Hom}, \texttt{\textbackslash RHom}), with extra wrapper packages (bm, bbm, dsfont, xspace, cancel, mathtools). The delivered numbering is the unit's own. Assets: no retained span contains a figure environment, a graphics-inclusion command, a PDF-page inclusion, a table or a TikZ picture; no image file is copied into this package, the package declares no asset dependency and no image file is redistributed with it. Licence: version arXiv:2309.02372v4 of this article is distributed under the Creative Commons Attribution 4.0 International licence, as recorded in the arXiv licence metadata for this exact version and shown on the abstract page https://arxiv.org/abs/2309.02372v4. A full rights notice is delivered at licenses/arxiv-2309.02372-CC-BY-4.0.txt. Characters: every retained excerpt is pure ASCII, so the delivered engine, XeLaTeX with its default Latin Modern text font, reports no missing character.
\begin{theorem}\label{observation}
Let $\varphi\colon R\rightarrow A$ be a ring homomorphism, where $R$ is a commutative noetherian ring and $A$ is a finite $R$-algebra.
The following two conditions are equivalent$\colon$
\begin{enumerate}
\item\label{basic}  $\Gdim_R(A)<\infty$ and $\RHom_R(A,R)$ is perfect over both $A$ and $A^{\rm op}$.

\item\label{TF} $\varphi$ has ascent and descent of finite Gorenstein dimension property.
\end{enumerate}
\end{theorem}

\begin{lemma}\label{TFtest}
Let $A$ be a noetherian ring. If $M$ is a finitely generated  $A$-module with finite Gorenstein dimension and $\Ext^i_A(M,A)=0$ for all $i>0$, then $M$ is Gorenstein projective.
\end{lemma}

\begin{theorem}\label{ADTF}
Let $\varphi\colon R\rightarrow A$ be a ring homomorphism, where $R$ is a commutative noetherian ring and $A$ is a finite $R$-algebra.
The following two conditions are equivalent$\colon$
\begin{enumerate}
\item  $A$ is Gorenstein projective over $R$ and $\Hom_R(A,R)$ is projective over both $A$ and $A^{\rm op}$.

\item $\varphi$ has ascent and descent of Gorenstein projective property.
\end{enumerate}
Moreover, if $\varphi$ has ascent and descent of Gorenstein projective property, then for each finitely generated $A$-module $M$, $\Gdim_R(M)=\Gdim_A(M)$.
\end{theorem}

\begin{proof}
The equality $\Gdim_R(M)=\Gdim_A(M)$ is straightforward under the given assumption. We only need to prove the first statement.

$(1)\Rightarrow (2)$. It follows from Theorem \ref{observation} that $\varphi$ has ascent and descent of finite Gorenstein dimension property. It suffices to prove that for any finitely generated $A$-module $M$, $\Gdim_A(M)=0$ if and only if $\Gdim_R(M)=0$.
By an analogous argument, the assertion for finitely generated $A^{\rm op}$-modules also holds.

Assume $\Gdim_A(M)=0$. Since $\varphi$ has descent of finite Gorenstein dimension property, $\Gdim_R(M)<\infty$. Moreover, by the hypothesis that $\Hom_R(A,R)$ is a projective  $A$-module, we infer that for all $i>0$,
$$\Ext_R^i(M,R)\cong \Ext_A^i(M,\Hom_R(A,R)) = 0.$$
It follows immediately from Lemma \ref{TFtest} that $\Gdim_R(M)=0$. 


Now, assume $\Gdim_R(M)=0$. Since $\Hom_R(A,R)$ is projective over $A^{\rm op}$ and $M$ is a Gorenstein projective $R$-module, we infer that $\Hom_R(A,R)\otimes_A M$ is also a Gorenstein projective $R$-module. Moreover, for any $i>0$ we conclude that
\begin{align*}
  \Ext^i_A(M, A)&\cong \Ext^i_A(M,\Hom_R(\Hom_R(A,R),R))\\
    &\cong \Ext_R^i(\Hom_R(A,R)\otimes_A M, R) = 0
\end{align*}
Note that $\Gdim_A(M)<\infty$ as $\varphi$ has ascent of finite Gorenstein dimension property. Then, it follows from Lemma \ref{TFtest} that $\Gdim_A(M)=0$. 
%Hence, $\varphi$ has ascent of Gorenstein projective property.


$(2)\Rightarrow (1)$. Since $\varphi$ has descent of Gorenstein projective property, $A$ is Gorenstein projective as an $R$-module. This yields that $\Hom_R(A,R)$ is also a Gorenstein projective $R$-module, and then $\Hom_R(A,R)$ is Gorenstein projective over both $A$ and $A^{\rm op}$ as $\varphi$ has ascent of Gorenstein projective property. There is a short exact sequence
$$0\rightarrow \Hom_R(A,R)\rightarrow F\rightarrow C\rightarrow 0$$
in $A\mbox{-mod}$, where $F$ is projective and $C$ is Gorenstein projective. Then the hypothesis yields that $C$ is also Gorenstein projective over $R$. We infer that the sequence is split from
$$\Ext^1_A(C,\Hom_R(A,R)) \cong  \Ext^1_R(C,R) = 0;$$
the isomorphism here follows from $\RHom_R(A,R)\simeq \Hom_R(A,R)$ and the adjunction $({\rm Res}, \RHom_R(A,-))$.
Hence, $\Hom_R(A,R)$ is a projective  $A$-module. The same argument will show that $\Hom_R(A,R)$ is a projective $A^{\rm op}$-module.
This completes the proof.
\end{proof}

\end{document}
